% ===== Block B13: 原文 p.116–130（p.130 为空白页，正文止于 p.129）| §5.6 抛物丛的模空间与 Hitchin 联络 | 译手 tr-ch5b =====

\subsection{抛物丛的模空间与 Hitchin 联络}

% [原文 p.116]

\noindent\textbf{5.6.1 抛物丛.}\quad 下面我们考察共形块空间与抛物丛模空间上某些线丛的整体截面空间之间的关系。


设 $G$ 为复数域 $\C$ 上的单连通单线性代数群，它对应于复单李代数 $\mathfrak g$。设 $T$ 为对应于 $\mathfrak g$ 的嘉当子代数 $\mathfrak h$ 的嘉当子群，$B$ 为对应于 Borel 子代数
\[
 \mathfrak b=\mathfrak h\oplus\Bigl(\bigoplus_{\alpha\in\Delta_+}\mathfrak g_\alpha\Bigr)
\]
的 Borel 子群。设 $\Sigma$ 是单根集 $\Pi$ 的一个子集，令
\[
 \mathfrak p_\Sigma=\mathfrak b\oplus\Bigl(\bigoplus_{\alpha\in\Sigma}\mathfrak g_{-\alpha}\Bigr),
\]
并设 $P_\Sigma$ 为 $G$ 中对应的闭子群。子群 $P_\Sigma$ 称为对应于子集 $\Sigma\subset\Pi$ 的\emph{标准抛物子群}。例如 $P_\varnothing=B$，而 $P_\Pi=G$。若权 $\lambda\in P_+$ 对所有 $\alpha\in\Sigma$ 满足 $\lambda(H_\alpha)=0$，则由指数映射它定义了 $P_\Sigma$ 的一个特征标。反之，任一整支配权 $\lambda$ 都决定一个标准抛物子群 $P_\Sigma$，其中
\[
 \Sigma=\{\alpha\in\Pi\mid\lambda(H_\alpha)=0\}.
\]


现在考虑亏格为 $g$ 的带坐标的 $N$-点黎曼曲面
\[
 \cX=(C;Q_1,\dots,Q_N;\xi_1,\dots,\xi_N).
\]
以下我们使用概形论的语言，即把 $C$ 视为 $\C$ 上的光滑曲线。固定一点 $P\in C\setminus\{Q_1,\dots,Q_N\}$ 及其坐标 $\xi$，令
\begin{align}
 \cX_0&=(C;P,Q_1,\dots,Q_N;\xi,\xi_1,\dots,\xi_N), \tag{5.40}\\
 C^*&=C\setminus\{P\}. \nonumber
\end{align}

对任一 $\C$-代数 $R$，我们对应一个群 $G(R((\xi)))$，其中 $R((\xi))$ 是系数取自 $R$ 的形式 Laurent 级数环。由此定义了从 $\C$-代数范畴到群范畴的一个群函子 $LG$。

为了在 $LG$ 上引入 ind-概形结构（即由概形的归纳极限给出的结构），我们选取一个忠实表示 $G\subset SL_r$ 并在以下固定之。对任一 $\C$-代数 $R$ 及任一整数 $m\ge 0$，令
\begin{align}
 LG(m)(R)=\bigl\{&A(\xi)\in G(R((\xi)))\subset SL_r(R((\xi)))\mid\text{ $A(\xi)$} \nonumber\\
 &\text{与 $A(\xi)^{-1}$ 的所有系数具有阶数 $\le N$ 的极点}\bigr\}. \nonumber
\end{align}

% [Tutor Note: 上面 $LG(m)$ 定义里的极点阶数界 $\le N$ 与参数 $m$ 不一致，应为 $\le m$：
%  紧接着的 $L^+G=LG(0)$ 只有在界为 $m$ 时才是"系数全为 $\xi$ 的非负幂"的群概形。
%  照原文排，见 ERRATA。]

于是 $LG(m)$ 是从 $\C$-代数范畴到集合范畴的函子，它可由一个无穷维仿射 $\C$-概形表示。特别地，
\[
 L^+G=LG(0)
\]
是仿射群概形。$\C$-群函子 $LG$ 视为一个 ind-群，即群概形序列 $\{LG(m)\}_{m\ge 0}$ 的归纳极限。这个 ind-群 $LG$ 称为 $G$ 的\emph{环路群}（loop group）。


% [原文 p.117]

对 $\C$-代数 $R$，令
\[
 A_{C_R}=H^0(C_R^*,\cO^*_{C_R}),
\]
其中 $C_R=C\times\operatorname{Spec}R$，$C_R^*=C^*\times\operatorname{Spec}R$。于是 $R\mapsto G(A_{C_R})$ 定义了从 $\C$-代数范畴到群范畴的函子 $L_CG$。容易证明 $L_CG$ 具有由 $LG$ 的结构诱导出的 ind-群结构。


$\mathfrak g$ 的\emph{环路代数}（loop algebra）$L\mathfrak g$ 定义为
\begin{equation}
 L\mathfrak g=\mathfrak g\otimes_\C\C((\xi)). \tag{5.41}
\end{equation}
环路代数 $L\mathfrak g$ 可以看作 $G$ 的环路群 $LG$ 的李代数：这是因为对任一 $\C$-代数 $R$，同态 $LG(R[[\eps]])\to LG(R)$ 的核等于 $L\mathfrak g(R)=\mathfrak g\otimes_\C R((\xi))$，其中 $\eps$ 是对偶数，即 $\eps^2=0$。类似地，$L_CG$ 的李代数 $L_C\mathfrak g$ 由 $\mathfrak g\otimes A_C$ 给出，即
\begin{equation}
 L_C\mathfrak g=\mathfrak g\otimes_\C H^0(C,\cO_C(*Q)). \tag{5.42}
\end{equation}


环路代数 $L\mathfrak g$ 的中心扩张是仿射李代数 $\hat{\mathfrak g}$。环路群通过伴随作用作用在 $L\mathfrak g$ 上，且该作用可扩张为 $\hat{\mathfrak g}$ 上的作用：
\[
 \mathrm{Ad}(\gamma)(X\otimes f(\xi),s\cdot c)=\Bigl(\mathrm{Ad}(\gamma)(X\otimes f(\xi)),\,
 s\cdot c+\operatorname*{Res}_{\xi=0}\Bigl(\gamma^{-1}\frac{d\gamma}{d\xi},
 (X\otimes f(\xi))\,d\xi\Bigr)\cdot c\Bigr).
\]

% [Tutor Note: 上式左端"$=\gamma(\cdots)$"处原文在等号后印了一个残缺的小字形
%  （形似 $\gamma$ 的断笔），按意义只能是右端的左括号，故照结构排为圆括号，
%  见 ERRATA。此式即中心扩张上的余循环：第二项是标准的 2-余循环
%  $\frac{1}{2\pi\ii}\int_\Gamma(\gamma^{-1}\dd\gamma\wedge X\otimes f)$ 的留数写法；
%  核对归一化：取 $X\otimes f=\gamma^{-1}\dd\gamma/\dd\xi$ 之外的单位元附近，
%  $\gamma=\ee^{tX}$ 时第二项给出 $t\,\mathrm{Res}(X,X\,d\xi)$，与
%  $\mathrm{Ad}$ 的线性化一致。]


其中 $(\ ,\ )$ 是施加于系数上的归一化 Cartan-Killing 型。

让我们考虑水平为 1 的左 $\hat{\mathfrak g}$-模 $\cH_0$ 以及中心扩张
\begin{equation}
 1\to\mathbb G_m\to GL(\cH_0)\to PGL(\cH_0)\to 1. \tag{5.43}
\end{equation}


$\hat{\mathfrak g}$ 在 $\cH_0$ 上的作用定义了 $\hat{\mathfrak g}$ 在 $PGL(\cH_0)$ 上的射影表示，而这个射影表示可以提升为 $LG$ 在 $PGL(\cH_0)$ 上的表示；再把正合列 (5.43) 沿 $LG$ 拉回，便得到 $LG$ 的一个中心扩张：
\begin{equation}
 1\to\mathbb G_m\to\widehat{LG}\to LG\to 1. \tag{5.44}
\end{equation}
（见 [LS]。）容易看出 $\hat{\mathfrak g}$ 是 $\widehat{LG}$ 的李代数。


注意到最高权向量 $|0\rangle\in\cH_0$ 生成直线 $\C|0\rangle$，它在 $\widehat{L^+G}$ 的作用下稳定。因此它定义了特征标
\begin{equation}
 \chi\colon\widehat{L^+G}\to\mathbb G_m, \tag{5.45}
\end{equation}
它给出 (5.44) 在 $\widehat{L^+G}$ 上的分裂。由于 $LG$ 的每个特征标都是平凡的，这个分裂是唯一的。


Grassmann 函子 $Q_G$ 定义为这样的函子：它对 $\C$-代数 $R$ 对应集合 $Q_G(R)$，其元素为对 $(E,\rho)$，其中 $E$ 是 $C$ 上的 $G$-丛，$\rho$ 是 $E$ 在 $C_R^*=(C\setminus\{P\})\times R$ 上的平凡化。函子 $Q_G$ 可由 $\C$-空间 $\mathcal Q_G$ 表示，且同构于 $LG/L^+G=\widehat{LG}/\widehat{L^+G}$（见 [LS]）。

% [Tutor Note: 原文此处印作 "of paris $(E,\rho)$"，应为 "pairs"（对），见 ERRATA。
%  $\mathcal Q_G$ 就是仿射 Grassmann 流形；集合 $LG/L^+G$ 与"$G$-丛 + $C^*$ 上平凡化"
%  的对应是 Beauville-Laszlo / 仿射 Grassmann 的标准描述：用 $\rho$ 把 $E$ 在
%  $C^*$ 上平凡化，再沿 $P$ 的形式邻域粘贴，转移函数即 $LG$ 的元素。]


设 $\mathcal M_G$ 为 $C$ 上 $G$-丛的栈，即这样的函子：它对 $\C$-代数 $R$ 对应 $C$ 上 $G$-丛的集合 $\mathcal M_G(R)$，使得 $\mathcal M_G(R)$ 的态射是 $G$-丛的同构。于是我们有忘记函子
\begin{align}
 \mathcal Q_G&\to\mathcal M_G \nonumber\\
 (E,\rho)&\mapsto E. \nonumber
\end{align}

在等同 $\mathcal Q_G=\widehat{LG}/\widehat{L^+G}$ 之下，(5.45) 的特征标 $\chi$ 定义了 $\mathcal Q_G$ 上的线丛 $\mathcal L$，它是 $\widehat{LG}$-线化的，从而在 $\mathcal M_G$ 上定义了一个线丛，

% [原文 p.118]

（仍记以同一符号 $\mathcal L$。）注意 $\mathcal L$ 的一个截面是 $\widehat{LG}$ 上的函数 $f$，满足
\[
 f(gh)=\chi(h)f(g),\qquad g\in\widehat{LG}(R),\ h\in\widehat{L^+G}(R).
\]
已知 $\mathcal L$ 是 $\operatorname{Pic}\mathcal Q_G$ 的一个正生成元。


借助等同 $\mathcal Q_G=LG/L^+G$，代数栈 $\mathcal M_G$ 典范地同构于双商 $L_\cX G\backslash LG/L^+G$，而投影
\[
 \mathcal Q_G\xrightarrow{\ \pi\ }\mathcal M_G
\]
与上面给出的忘记态射相同，并且它对 étale 拓扑是局部平凡的（[LS]）。自然投影诱导拉回态射
\[
 \pi^*\colon\operatorname{Pic}(\mathcal M_G)\to\operatorname{Pic}(\mathcal Q_G).
\]
Laszlo 与 Sörger [LS] 证明 $\pi^*$ 是同构，从而 $\operatorname{Pic}(\mathcal M_G)$ 由 $\mathcal M_G$ 上的线丛 $\mathcal L$ 生成。


现在定义抛物 $G$-丛。固定 $\vec\lambda=(\lambda_1,\dots,\lambda_N)\in P_\ell^N$。于是每个 $\lambda_j$ 决定一个子集
\[
 \Sigma_j=\{\alpha\in\Delta_+\mid\lambda_j(H_\alpha)=0\}
\]
以及对应的标准抛物子群 $P_j=P_{\Sigma_j}$。此外，$\lambda_j$ 通过指数映射决定子群 $P_j$ 的一个特征标 $p_j$。再令 $\Gamma_j=\Pi\setminus\Sigma_j$，并考虑正整数 $m_{j,k}=\lambda_j(H_{\alpha_k})$，其中 $\alpha_k\in\Gamma_j$。设 $E$ 是一个 $G$-丛。由于群 $G$ 作用在 $E$ 上，我们可以定义相伴的 $G/P_j$-丛 $E(G/P_j)$。



\begin{definition}[5.24]
设 $\vec P=(P_1,\dots,P_N)$ 为附着于点集 $(Q_1,\dots,Q_N)$ 的一族标准抛物子群 $P_j=P_{\Sigma_j}$，并设 $\underline m=(\vec m_1,\dots,\vec m_N)$ 为一族严格正整数，其中 $\vec m_j=(m_{j,k})_{\alpha_k\in\Gamma_j}$，$\Gamma_j=\Pi\setminus\Sigma_j$。$\cX$ 上型为 $\vec P$ 的\emph{拟抛物 $G$-丛}是指曲线 $C$ 上的一个 $G$-丛，并附带对每个 $1\le j\le N$ 选取的元素 $F_j\in E(G/P_j)(Q_j)$。$\cX$ 上型为 $(\vec P,\underline m)$ 的\emph{抛物 $G$-丛}是指 $\cX$ 上的一个拟抛物 $G$-丛 $E$，并附带对每个 $1\le j\le N$ 的抛物权 $(m_{j,k})_{\alpha_k\in\Gamma_j}$，其中 $m_{j,k}$ 是正整数且 $\underline m=(m_{j,k})$。
\end{definition}


对 $\C$-代数 $R$，由 $S=\operatorname{Spec}R$ 参数化的型为 $\vec P$ 的拟抛物 $G$-丛族 $E$ 是指 $C_R=C\times_\C R$ 上的一个 $G$-丛，并附带 $N$ 个截面
\[
 \sigma_j\colon S\to E(G/P_i)|_{S\times\{Q_j\}}.
\]
从 $(E,\vec\sigma)$ 到 $(E',\vec\sigma')$ 的态射是指 $G$-丛的态射 $f\colon E\to E'$，使得 $\sigma'_j=f|_{S\times\{Q_j\}}\circ\sigma_j$。


我们可以定义一个从 $\C$-代数范畴到群胚范畴的伪函子：对 $\C$-代数 $R$ 对应以"由 $\operatorname{Spec}R$ 参数化的型为 $\vec P$ 的拟抛物 $G$-丛族"为对象、以族之间的同构为箭头的群胚。对任一同态 $R\to R'$，我们在相应的群胚之间有自然的函子。由此定义了型为 $\vec P$ 的拟抛物 $G$-丛的 $\C$-栈 $\mathcal M_G^{pra}(\cX,\vec P)$。

对任一环同态 $R\to R'$，我们可以定义型为 $\vec P$ 的拟抛物 $G$-丛的 $\C$-栈 $\mathcal M_G^{par}(\cX,\vec P)$，它是相应群胚之间的函子。

% [Tutor Note: 上面两段在原文里几乎是重复的：前一段用 $\mathcal M_G^{pra}$（"pseudo"
%  之缩写）定义了拟抛物丛栈，后一段又用 $\mathcal M_G^{par}$ 定义"拟抛物"丛栈，
%  只把 $R\to R'$ 的函子性重述一遍。按上下文后一段应是在定义"抛物"丛（带权）的
%  栈，原文很可能把 parabolic 写成了 quasi-parabolic；照原文译，见 ERRATA。
%  核对：下一段的 $\mathcal Q_G^{par}$ 用 $\prod_j G/P_j$ 只记录拟抛物数据，
%  故 $\mathcal Q_G^{par}$ 只表示 $\mathcal M_G^{pra}$ 的覆盖，不含权。]
%  另：$N$ 个截面的定义式原文印 $\sigma_j\colon S\to E(G/P_i)|_{S\times\{Q_j\}}$；
%  按 $\sigma_j$ 的取值应为 $G/P_j$，译文照排原文（见 ERRATA）。

现在定义
\[
 \mathcal Q_G^{par}(\cX,\vec P)=\mathcal Q_G\times\prod_{j=1}^N G/P_j.
\]


% [原文 p.119]

于是 ind-群 $L_CG$ 作用在 $\mathcal Q_G$ 上；并且在 $Q_j$ 处通过赋值 $\mathrm{ev}(Q_j)\colon L_CG\to G\to G/P_j$，我们得到 $L_CG$ 在 $\mathcal Q_G^{par}(\cX,\vec P)$ 上的自然作用。于是可以证明下述定理。

\begin{theorem}[5.25] (Laszlo \& Sörger [LS])
存在栈的典范同构
\[
 \bar\pi\colon L_CG\backslash\mathcal Q_G^{par}(\cX,\vec P)\simeq\mathcal M_G^{par}(\cX,\vec P).
\]
投影映射 $\pi\colon\mathcal Q_G^{par}(\cX,\vec P)\to\mathcal M_G^{par}(\cX,\vec P)$ 对 étale 拓扑是局部平凡的。
\end{theorem}


另外，Laszlo 与 Sörger [LS] 还证明了以下结果。

\begin{theorem}[5.26]
\begin{equation}
 \operatorname{Pic}\bigl(\mathcal Q_G^{par}(\cX,\vec P)\bigr)=\Z\mathcal L\times\prod_{j=1}^N\operatorname{Pic}(G/P_j)
 =\Z\mathcal L\times\prod_{j=1}^n X(P_j). \tag{5.46}
\end{equation}
其中 $X(P_j)$ 是 $P_j=P_{\Sigma_j}$ 的特征标群。
\end{theorem}

% [Tutor Note: (5.46) 右端第二个乘积的上限印作小写 $n$，而第一个乘积是 $N$；
%  两处必须相同，应为 $N$（$N$ 是标记点个数，$n$ 全书未作此用）。照原文排，
%  见 ERRATA。另：$\operatorname{Pic}(G/P_j)$ 是自由 $\Z$-模，秩等于 $P_j$ 的
%  半单秩（即 $\Gamma_j$ 的基数），与 $X(P_j)$ 同构；故两端一致。]

设 $\vec\lambda=(\lambda_1,\dots,\lambda_N)$ 为附着于点集 $(Q_1,\dots,Q_N)$ 的标号。令
\[
 \Sigma_j=\{\alpha\in\Pi\mid\lambda_j(H_\alpha)=0\},\qquad \Gamma_j=\Pi\setminus\Sigma_j,
\]
并再令
\begin{align*}
 \vec P&=(P_1,\dots,P_N),\qquad P_j=P_{\Sigma_j},\\
 \underline m&=(m_{j,k}),\qquad m_{j,k}=\lambda_j(H_{\alpha_k}),\quad \alpha_k\in\Gamma_j.
\end{align*}


设 $\mathcal L$ 为上文中在 $\mathcal Q_G$ 与 $\mathcal M_G$ 上定义过的线丛。对任一非负整数 $\ell$，定义线丛
\[
 \mathcal L(\ell,\underline m)=\mathcal L^\ell\otimes\Bigl(\bigotimes_{j=1}^N\bigl(\otimes_{\alpha_k\in\Gamma_j}\Z\mathcal L_{j,k}^{m_{j,k}}\bigr)\Bigr)
\]
（在栈 $\mathcal M_G^{par}(\cX,\vec P)$ 上），其中 $\mathcal L_{j,k}$ 是 $G/P_j$ 上由与 $\alpha_k\in\Gamma_j$ 相伴的特征标所对应的线丛。于是 $\mathcal L(\ell,\underline m)$ 到 $\mathcal Q_G^{par}(\cX,\vec P)$ 的拉回由
\[
 \pi^*\mathcal L(\ell,\underline m)=\cO_{\mathcal Q_G}(\ell)\otimes\Bigl(\bigotimes_{j=1}^N\mathcal L_{-\lambda_j}\Bigr)
\]
给出，其中 $\mathcal L_{-\lambda_j}$ 是齐性空间 $G/P_j$ 上由对应于权 $-\lambda_j$ 的特征标所定义的线丛。

% [Tutor Note: 公式里 $\Z\mathcal L_{j,k}^{m_{j,k}}$ 的因子 $\Z$ 在全书中别处未出现，
%  按张量积结构此处应是线丛 $\mathcal L_{j,k}^{m_{j,k}}$ 本身（$\Z$ 很可能是笔误或
%  排版残留）。照原文排，见 ERRATA。另：拉回式里 $\mathcal L^\ell$ 变成
%  $\cO_{\mathcal Q_G}(\ell)$、$\mathcal L_{j,k}^{m_{j,k}}$ 变成 $\mathcal L_{-\lambda_j}$，
%  关键是 $\sum_k m_{j,k}\alpha_k=\lambda_j$（因为 $\lambda_j=\sum_k\lambda_j(H_{\alpha_k})\varpi_k$
%  且 $\alpha_k$ 与 $\varpi_k$ 对偶），故张量积按权相加恰得 $-\lambda_j$ 的线丛。]

\begin{theorem}[5.27] (Laszlo \& Sörger [LS])
存在典范同构
\[
 \mathcal V_{\vec\lambda}^\dagger(\cX)\simeq H^0\bigl(\mathcal M_G^{par}(\cX,\vec P),\mathcal L(\ell,\underline m)\bigr).
\]
\end{theorem}


\noindent\textit{证明概述.}\quad 首先注意，由定理 3.15 与定理 3.18，若令
\[
 \cX_0=(C;P,Q_1,\dots,Q_N;\xi,\xi_1,\dots,\xi_N),
\]
则有典范同构
\[
 \mathcal V_{0\vec\lambda}^\dagger(\cX_0)\simeq\mathcal V_{\vec\lambda}^\dagger(\cX)
 \simeq\bigl\{\langle\Phi|\in\cH_0^\dagger\otimes V_{\vec\lambda}^*\mid
 \langle\Phi|\mathfrak g\otimes H^0(C,\cO_C(*P))\bigr\}.
\]
由 Borel-Bott 定理，有右 $\mathfrak g$-模同构
\[
 H^0(G/P_j,\mathcal L_{-\lambda_j})\simeq V_{\lambda_j}^*.
\]

% [Tutor Note: 上式（集合描述）里最后一个条件漏印了 "$=0$"：应为
%  $\langle\Phi|(\mathfrak g\otimes H^0(C,\cO_C(*P)))=0$，即"被 $\mathfrak g\otimes H^0$
%  零化"，这正是定理 3.15/3.18 给出的共形块描述。照原文排，见 ERRATA。
%  核对：若不加 $=0$，右端是整个 $\cH_0^\dagger\otimes V^*_{\vec\lambda}$，
%  维数无穷，与左边有限维的 $\mathcal V^\dagger_{\vec\lambda}(\cX)$ 不可能同构。]


% [原文 p.120]

此外，已知
\begin{equation}
 H^0(\mathcal Q_G,\cO_{\mathcal Q}(\ell))\simeq\cH_0^\dagger \tag{5.47}
\end{equation}
其中 $\cH_0^\dagger$ 是水平为 $\ell$、最高权为 0 的不可约可积最高权 $\hat{\mathfrak g}$-模。（见 [Ku]、[Ma]。由 $\mathcal Q_G$ 的定义，$H^0(\mathcal Q_G,\cO_{\mathcal Q}(\ell))$ 具有自然的 $\hat{\mathfrak g}$-模结构。若把 $\cH_0^\dagger$ 视为 §2.2.1 中定义左 $\hat{\mathfrak g}$-模那样，则同构 (5.47) 是左 $\hat{\mathfrak g}$-模同构。）因此有 $\widehat{LG}$-模的典范同构
\[
 H^0(\mathcal Q_G,\cO_{\mathcal Q}(\ell))\otimes\Bigl(\bigotimes_{j=1}^N H^0(G/P_j,\mathcal L_{-\lambda_j})\Bigr)
 \simeq\cH_0^\dagger\otimes V_{\vec\lambda}^\dagger.
\]


另一方面，我们有 $\hat{\mathfrak g}$-模的同构：
\[
 H^0\bigl(\mathcal Q_G^{par}(\cX',\vec P),\pi^*(\mathcal L(\ell,\underline m))\bigr)
 \simeq H^0(\mathcal Q_G,\cO_{\mathcal Q}(\ell))\otimes\Bigl(\bigotimes_{j=1}^N H^0(G/P_j,\mathcal L_{-\lambda_j})\Bigr).
\]
此外还存在同构
\[
 \Bigl[H^0\bigl(\mathcal Q_G^{par}(\cX',\vec P),\pi^*(\mathcal L(\ell,\underline m))\bigr)\Bigr]^{L_CG}
 \simeq\Bigl[H^0\bigl(\mathcal Q_G^{par}(\cX',\vec P),\pi^*(\mathcal L(\ell,\underline m))\bigr)\Bigr]^{L_C\mathfrak g},
\]
其中 $L_C\mathfrak g$ 是 $L_CG$ 的李代数（见 (5.42)）。最后，可以证明下述同构：
\[
 H^0\bigl(\mathcal M_G^{par}(\cX',\vec P),\mathcal L(\ell,\underline m)\bigr)
 \simeq\Bigl[H^0\bigl(\mathcal Q_G^{par}(\cX',\vec P),\pi^*(\mathcal L(\ell,\underline m))\bigr)\Bigr]^{L_CG}.
\]
这就证明了该定理。\hfill$\square$


\noindent\textbf{5.6.2 Hitchin 联络.}\quad 本小节我们引入 Hitchin 联络，并说明它与我们的联络本质上是同一个。

先对带点光滑曲线族叙述定理 5.27 的相对版本。设 $\mathfrak F=(\pi\colon\mathcal C\to S;s)$ 是仿射概形 $S=\operatorname{Spec}(R)$ 上非常值的带点光滑曲线族。我们可以定义 $S$ 上的 ind-群 $L_\pi G$ 与 $L_\pi^+G$，作为 $LG$ 与 $L^+G$ 的相对版本。同样，对族 $\mathfrak F$，也可以像 §5.6.1 那样定义 $S$ 上的 ind-群概形 $\widehat{L_\pi G}$（相应地 $\widehat{L_\pi^+G}$），它们是 $\widehat{LG}$（相应地 $\widehat{L^+G}$）的相对版本。此外，我们可以把 ind-群族 $(L_{C_x}G)_{x\in S}$（$C_x=\pi^{-1}(x)$）粘起来，定义 $S$ 上的 ind-群 $L_{\mathfrak F}G$。存在嵌入 $L_{\mathfrak F}G\hookrightarrow\widehat{L_\pi G}$。与 (5.45) 类似，我们可以定义 $S$ 上的特征标
\begin{equation}
 \chi\colon\widehat{L_\pi^+G}\to\mathbb G_{m,S}. \tag{5.48}
\end{equation}


现在令
\[
 \mathcal Q_\pi=\widehat{L_\pi G}/\widehat{L_\pi^+G}.
\]
这是 §5.6.1 中 $\mathcal Q_G$ 的相对版本，并且存在自然态射
\[
 q_\pi\colon\mathcal Q_\pi\to S.
\]
态射 $q_\pi$ 在 Zariski 拓扑下局部平凡，因为沿截面 $s$ 选取一个形式坐标就给出这样的平凡化。又，特征标 $\chi^{-1}\colon\widehat{L_\pi^+G}\to\mathbb G_{m,S}$ 定义了 $\mathcal Q_\pi$ 上的线丛 $\mathcal L$，其中 $\chi$ 由 (5.48) 给出。由于 $L_{\mathfrak F}G$ 嵌入 $\widehat{L_\pi G}$，线丛 $\mathcal L$ 是 $L_{\mathfrak F}G$-线化的，从而它

定义了商栈 $\mathcal M_{G,\pi}=L_{\mathfrak F}G\backslash\mathcal Q_\pi$ 上的一个线丛 $\mathcal L$，这个商栈是 $G$-丛的模栈的族。

% [原文 p.121]

现在假定沿截面 $s$ 给定了一个形式坐标 $\xi$。这时我们可以定义同构
\[
 \iota_\pi\colon q_{\pi*}(\mathcal L^{\otimes\ell})\simeq\cH_0^\dagger\otimes\cO_S,
\]
其中 $\cH_0^\dagger$ 是水平为 $\ell$、最高权为 0 的可积最高权右 $\hat{\mathfrak g}$-模。这是同构 (5.47) 的相对版本。此外，该同构还诱导出同构
\begin{equation}
 \mathcal V_0(\mathfrak F)^\dagger\simeq p_{\pi*}(\mathcal L^{\otimes\ell}), \tag{5.49}
\end{equation}
其中 $p_\pi\colon\mathcal M_{G,\pi}=L_{\mathfrak F}G\backslash\mathcal Q_\pi\to S$ 是自然态射，因为我们可以证明同构
\[
 p_{\pi*}(\mathcal L^{\otimes\ell})=\bigl(q_{\pi*}(\mathcal L^{\otimes\ell})\bigr)^{L_{\mathfrak F}G}.
\]

注意，上述论证对态射 $\pi$ 是函子性的。设 $\mathcal M_{g,1}$（相应地 $\mathcal M_g$）是亏格 $g$ 的光滑曲线（相应地带一个点的光滑曲线）的模栈，$\pi\colon\cX\to\mathcal M_{g,1}$ 是亏格 $g$ 的通用曲线，带标准截面 $s$。设 $\mathcal M_G$ 是光滑亏格 $g$ 曲线上半稳定 $G$-丛的通用模栈，$\underline{\mathcal M}_G$ 是光滑亏格 $g$ 曲线上半稳定 $G$-丛的通用粗模空间。令 $p\colon\mathcal M_G\to\mathcal M_g$ 是自然态射。于是我们可以应用上述论证，从而存在 $\mathcal M_G$ 上的线丛 $\mathcal L$。


又注意，由真空传播（§3.3），共形块 $\mathcal V_0((C;P;\xi))$ 与 $\mathcal V_0((C;Q;\eta))$ 在标量乘法意义下同构。因此射影丛
\[
 \mathbf{P}\mathcal V_0^\dagger:=\Bigl(\mathcal V_0^\dagger\bigl((\cX\to\mathcal M_{g,1};s)\bigr)\setminus\{\text{零截面}\}\Bigr)/\C^*
\]
下降到 $\mathcal M_g$ 上的一个射影丛，我们仍用同一记号 $\mathbf{P}\mathcal V_0^\dagger$ 表示它。于是，定理 5.27 的相对版本是说我们有射影丛的同构：
\begin{equation}
 \mathbf{P}\mathcal V_0^\dagger\simeq\mathbf{P}p_*(\mathcal L^{\otimes\ell}):=\Bigl(p_*(\mathcal L^{\otimes\ell})\setminus\{\text{零截面}\}\Bigr)/\C^*. \tag{5.50}
\end{equation}

Hitchin（[Hi]）在 $\mathbf{P}p_*(\mathcal L^{\otimes\ell})$ 上定义了一个联络。为定义 Hitchin 联络，让我们先考虑光滑射影簇、其上的复线丛以及该线丛的一个整体截面所成的三元组的形变理论。设 $X$ 是复数域 $\C$ 上的光滑射影簇，$L$ 是 $X$ 上的复线丛，$s\in H^0(X,L)$ 是一个截面。我们感兴趣的是对 $(X,L)$ 以及三元组 $(X,L,s)$ 的一阶形变作分类。

用 $(X_\eps,L_\eps,s_\eps)$ 记 $(X,L,s)$ 在 $\operatorname{Spec}\C[\eps]$（其中 $\eps^2=0$）上的一个形变。设 $\{U_i\}_{i\in I}$ 是 $X$ 的仿射开覆盖。由于 $U_i$ 是光滑且仿射的，$X_\eps$ 到 $U_i$ 上的限制 $U_{i,\eps}$ 是平凡的。我们选取一个同构
\[
 \iota_i\colon\cO_{U_i}[\eps]=\cO_{U_i}\otimes_\C\C[\eps]\simeq\cO_{U_{i,\eps}},
\]
使得到 $\eps=0$ 的限制是恒等。则 $\iota_j^{-1}\circ\iota_i$ 由一个矩阵
\[
 \begin{pmatrix} \id & 0\\ \xi_{ij} & \id \end{pmatrix}
\]
给出，其中 $\xi_{ij}$ 是 $\cO_{U_i\cap U_j}$ 的导子。而且 $\{\xi_{ij}\}$ 是一个上循环，给出 $H^1(X,\Theta_X)$ 中的一个元素，它就是形变 $X_\eps$ 的小平–Spencer 类（见 §1.1.3）。

% [原文 p.122]

另外，$L_\eps$ 到 $U_i$ 上的限制 $LU_{i,\eps}$ 也是平凡的。我们选取一个同构
\[
 \phi_i\colon LU_i[\eps]=LU_i\otimes_\C\C[\eps]\simeq LU_{i,\eps},
\]
使得到 $\eps=0$ 的限制是恒等。同构 $\phi_j^{-1}\circ\phi_i$ 由一个矩阵
\[
 \begin{pmatrix} \id & 0\\ \eta_{ij} & \id \end{pmatrix}
\]
给出，其中 $\eta_{ij}$ 是 $L_{U_i\cap U_j}$ 上的一个一阶微分算子，其象征为 $\eta_{ij}$。

% [Tutor Note: 这里的"其象征为 $\eta_{ij}$"是原文的笔误，应为"其象征为 $\xi_{ij}$"：
%  上一页的矩阵给出 $X_\eps$ 的粘接，其 $(2,1)$-元 $\xi_{ij}$ 是 $\cO$ 的导子；本页对
%  $L_\eps$ 的粘接，$(2,1)$-元 $\eta_{ij}$ 是一阶微分算子，取象征（即最高阶项）正好
%  回到 $\xi_{ij}$。若照"象征为 $\eta_{ij}$"去读，下文"取象征把 $\{\eta_{ij}\}$ 送到
%  $\{\xi_{ij}\}$"就无法成立。此处按原文排版，未改。]

令 $\mathcal D_X^i(L)$ 是作用在 $L$ 上的、阶数不超过 $i$ 的微分算子层。则 $\{\eta_{ij}\}$ 定义了 $H^1(X,\mathcal D_X^1(L))$ 中的一个元素，它就是形变 $(X_\eps,L_\eps)$ 的小平–Spencer 类。通过对微分算子取象征，元素 $\{\eta_{ij}\}\in H^1(X,\mathcal D_X^1(L))$ 被映到元素 $\{\xi_{ij}\}\in H^1(X,\Theta_X)$。

现在，截面 $s_\eps$ 到 $U_i$ 上的限制 $s_{U_{i,\eps}}$ 写成如下形式
\[
 s_{U_{i,\eps}}=\phi_i(s_{U_i}+\eps\sigma_i),
\]
其中 $\sigma_i\in H^0(U_i,L_{U_i})$ 被唯一确定。在 $U_i\cap U_j$ 上有 $s_{U_{i,\eps}}=s_{U_{j,\eps}}$。因此我们有等式
\begin{equation}
 \sigma_j-\sigma_i=\eta_{ij}(s). \tag{5.51}
\end{equation}
赋值映射 $ev_s\colon\mathcal D_X^i(L)\to L$ 由
\[
 ev_s(D)=D(s)
\]
定义。令 $d_1s$ 是复形
\[
 d_1s=\{\mathcal D_X^i(L)\xrightarrow{ev_s}L\},
\]
其中 $\mathcal D_X^i(L)$ 处于 0 次，$L$ 处于 1 次。则由 (5.51)，
\[
 (\eta_{ij},\sigma_i)\in C^1(\{U_i\},d_1s)=C^1(\{U_i\},\mathcal D_X^1(L))\oplus C^0(\{U_i\},L)
\]
是一个上循环，从而定义了该复形 $d_1s$ 的超上同调群 $\HH^1(d_1s)$ 中的一个元素。

% [Tutor Note: $d_1s$ 这个记号指"由截面 $s$ 决定的复形"，即两项复形
%  $\cD_X^1(L)\xrightarrow{ev_s}L$；它不是微分，容易与 $d_1s$ 读成"$d(s)$ 的一阶部分"。
%  相应地 $d_2s$ 是 $\cD_X^2(L)\xrightarrow{ev_s}L$。该复形的超上同调分类三元组的一阶
%  形变，这是 Hitchin 联络构造中标准的一步（原文未给证明）。
%  另：本页下文"三元组 $(X,L.s)$"中的句点是原文笔误，应为逗号 $(X,L,s)$；照排，见 ERRATA。]

{\emergencystretch=22.3pt 于是，三元组 $(X,L,s)$ 的一阶形变 $(X_\eps,L_\eps,s_\eps)$ 的同构类定义了复形 $d_1s$ 的上同调类 $[(\eta_{ij},\sigma_i)]\in\HH^1(d_1s)$。反过来，$\HH^1(d_1s)$ 的任一上同调类都决定三元组 $(X,L.s)$ 的一个一阶形变 $(X_\eps,L_\eps,s_\eps)$ 的同构类，并且这给出三元组 $(X,L,s)$ 的一阶形变的同构类集合与 $\HH^1(d_1s)$ 之间的一个同构。注意存在上同调群的同态\par}
\begin{equation}
 \beta\colon\HH^1(d_1s)\to H^1(X,\mathcal D_X^1(L)) \tag{5.52}
\end{equation}
它由附着于超上同调的谱序列给出。

现在，象征映射与赋值映射 $ev_s$ 给出下述交换图，它给出复形 $d_1s$、$d_2s$ 与 $\{S^2T_X\to0\}$ 的一个正合列：
\[
 \begin{array}{ccccccccc}
 0 & \longrightarrow & \mathcal D_X^1(L) & \longrightarrow & \mathcal D_X^2(L) & \longrightarrow & S^2T_X & \longrightarrow & 0\\[2mm]
 & & \big\downarrow\hbox{\scriptsize $ev_s$} & & \big\downarrow\hbox{\scriptsize $ev_s$} & & \big\downarrow & & \\[2mm]
 0 & \longrightarrow & L & \longrightarrow & L & \longrightarrow & 0 & \longrightarrow & 0
 \end{array}
\]
其中 $S^2T_X$ 是 $X$ 的全纯切丛 $T_X$ 的对称积。因此，由超上同调的正合列我们得到同态
\begin{equation}
 \delta_s\colon H^0(X,S^2T_X)\to\HH^1(d_1s). \tag{5.53}
\end{equation}

% [原文 p.123]

同态 $\delta_s$ 由下述方式给出。对 $w\in H^0(X,S^2T_X)$ 取 $D_i^{(2)}\in H^0(U_i,\mathcal D_X^2(L))$，使它在 $U_i$ 上被映到 $w|_{U_i}$。于是有
\[
 \delta_s(w)=(\{D_i^{(2)}-D_j^{(2)}\},\{D_i^{(2)}(s)\})\in\HH^1(d_1s).
\]
由 (5.52) 我们有
\[
 \beta(\delta_s(w))=\{D_i^{(2)}-D_j^{(2)}\}\in H^1(X,\mathcal D_X^1(L)),
\]
并且该上同调类与截面 $s$ 的选取无关。

% [Tutor Note: 末句"与截面 $s$ 的选取无关"按原文照译，但此处 $s$ 是一开始就固定
%  给定的，不可能再"选取"；从上下文看，应读作"与 $D_i^{(2)}$ 的选取无关"——
%  若换一组 $D_i^{(2)}$，两者之差处处是阶数不超过 1 的算子，故上同调类不变。]

现在让我们考虑 $L_\eps$ 上阶数至多为 2 的一个微分算子 $D$。在 $U_{i,\eps}$ 上，微分算子 $D$ 写成
\[
 D_i=\partial_\eps+R_i\colon LU_i[\eps]\to LU_i[\eps],
\]
其中
\[
 \partial_\eps(a+\eps b)=b,\qquad a,b\in LU_i,
\]
而 $R_i$ 是一个只涉及 $\operatorname{Spec}\C[\eps]$ 上纤维方向（即不涉及基方向 $\eps$）的导数的二阶微分算子。把 $D$ 的象征限制到 $X$ 上，记作 $\bar\rho(D)$。若写 $R_i=\bar R_i+\eps\bar R_i'$，则象征 $\bar\rho(D)\in H^0(X,S^2T_X)$ 由各 $\bar R_i$ 在 $U_i$ 上的象征 $\sigma(\bar R_i)$ 的汇集给出，其中 $\bar R_i,\bar R_i'\colon LU_i\to LU_i$ 是二阶微分算子。因此，超上同调类 $\delta_s(\bar\rho(D))$ 由
\[
 \delta_s(\bar\rho(D))=(\{\bar R_i-\bar R_j\},\{\bar R_i(s)\})
\]
给出。


于是，容易证明 $-\beta(\delta_s(\bar\rho(D)))$ 对应于一阶形变 $(X_\eps,L_\eps)$。又注意，$-\delta_s(\bar\rho(D))$ 对应于一阶形变 $(X_\eps,L_\eps,s_\eps)$，当且仅当
\[
 D_i(s_\eps)\equiv0\quad(\operatorname{mod}\eps).
\]

现在让我们考虑 $\C$ 上光滑射影簇的光滑族 $p\colon\cX\to S$。设 $\mathcal L$ 是 $\cX$ 上的复线丛。设 $\mathcal D_{\cX}^i(\mathcal L)$ 是 $\mathcal L$ 在 $\cX$ 上的阶数至多为 $i$ 的微分算子层。用 $\mathcal D_{\cX/S}^i(\mathcal L)\subset\mathcal D_{\cX}^i(\mathcal L)$ 记 $p^{-1}\cO_S$-线性算子的子层。对任意正整数 $i$，存在象征映射
\[
 \sigma^i\colon\mathcal D_{\cX}^i(\mathcal L)\to\operatorname{Sym}^i_{\cO_{\cX}}\mathcal T_{\cX},
\]
它把 $\mathcal D_{\cX/S}^i(\mathcal L)$ 映入 $\operatorname{Sym}^i_{\cO_{\cX}}\mathcal T_{\cX/S}$。下文我们用记号 $S^i\mathcal F$ 表示 $\operatorname{Sym}^i_{\cO_{\cX}}(\mathcal F)$。令
\[
 \mathcal W_{\cX/S}(\mathcal L)=\mathcal D_{\cX}^1(\mathcal L)+\mathcal D_{\cX/S}^2(\mathcal L)\subset\mathcal D_{\cX}^2(\mathcal L).
\]
层 $\mathcal W_{\cX/S}(\mathcal L)$ 是 $\mathcal L$ 上二阶微分算子的子层，其象征落在 $S^2\mathcal T_{\cX/S}$ 中。于是存在正合列：
\begin{equation}
 0\to\mathcal D_{\cX/S}^1(\mathcal L)\to\mathcal W_{\cX/S}(\mathcal L)\xrightarrow{\sigma}p^*(T_S)\oplus S^2\mathcal T_{\cX/S}\to0. \tag{5.54}
\end{equation}
令
\[
 \sigma_1=p_1\circ\sigma\colon\mathcal W_{\cX/S}(\mathcal L)\to p^*(T_S),
\]
其中 $\sigma$ 是 (5.54) 中的同态，而 $p_1\colon p^*(T_S)\oplus S^2\mathcal T_{\cX/S}\to p^*(T_S)$ 是到第一个因子的投影。

% [原文 p.124]

\begin{definition}[5.28]
$\mathcal L$ 上的一个热算子（heat operator）$D$ 是一个左 $\cO_{\cX}$-线性同态
\[
 D\colon p^*(T_S)\to\mathcal W_{\cX/S}\subset\mathcal D_{\cX}^2(\mathcal L),
\]
使得
\[
 \sigma_1\circ D=\id.
\]
\end{definition}

若取 $S$ 的坐标邻域 $U$ 与 $p^{-1}(U)\subset\cX$ 的邻域 $V$，其局部坐标分别为 $U$ 的 $(t_1,\dots,t_r)$ 与 $V$ 的 $(t_1,\dots,t_r,z_1,\dots,z_n)$，则热算子 $D$ 写成
\begin{equation}
 D\Bigl(\frac{\partial}{\partial t_k}\Bigr)=f_k+\frac{\partial}{\partial t_k}+\sum_{i=1}^n f_{k,i}\frac{\partial}{\partial z_i}+\sum_{1\le i,j\le n}f_{k,ij}\frac{\partial^2}{\partial z_i\partial z_j},\quad k=1,\dots,r \tag{5.55}
\end{equation}
其中 $f_k,f_{k,i},f_{k,ij}\in\cO_{\cX}(V)$。

% [Tutor Note: 热算子是本节的技术核心：它在"$\cX$ 上的二阶微分算子"与"$S$ 上的
%  联络"之间架桥。$D(\partial/\partial t_k)$ 的首项 $\partial/\partial t_k$ 保证
%  $\sigma_1\circ D=\id$（即基方向的一阶部分正好是 $\partial/\partial t_k$ 本身），
%  而 $\partial/\partial z_i$、$\partial^2/\partial z_i\partial z_j$ 各项给出沿纤维
%  方向的二阶修正；它们正是"热核"型算子的局部形状。]

$\mathcal L$ 在 $S$ 上的一个射影热算子（projective heat operator）$\bar D$ 是一个 $\cO_S$-线性映射
\[
 \bar D\colon T_S\to p_*\bigl(\mathcal W_{\cX/S}(\mathcal L)\bigr)/\cO_S,
\]
使得任一局部提升 $D_U\colon p^*(T_S)|_{p^{-1}(U)}\to\mathcal W_{\cX/S}(\mathcal L)|_{p^{-1}(U)}$ 都是 $\mathcal L$ 在 $U$ 上的热算子。

热算子 $D$ 如下地定义一个联络
\[
 \nabla(D)\colon p_*(\mathcal L)\to p_*(\mathcal L)\otimes_{\cO_S}\Omega^1_S.
\]
设有 $s\in p_*(\mathcal L)=H^0(p^{-1}(U),\mathcal L)$ 与 $\theta\in H^0(U,T_S)$。则 $D(p^{-1}\theta)$ 是 $\mathcal L$ 在 $p^{-1}(U)$ 上的二阶微分算子。因此我们有
\[
 D(p^{-1}\theta)(s)\in\mathcal L(p^{-1}(U))=p_*(\mathcal L)(U).
\]
此外，由 $D$ 的局部表达式 (5.55)，对任意 $f\in\cO_S(U)$ 我们有
\[
 D(p^{-1}T)(fs)=fD(p^{-1}T)(s)+\theta(f)s.
\]

% [Tutor Note: 上式中出现的 $T$ 是原文笔误，应为 $\theta$。理由是：$D(p^{-1}\theta)$
%  作用在 $fs$（$f$ 只依赖底空间 $S$）上，由 (5.55) 中 $\partial/\partial t_k$ 那一项
%  产生 $\theta(f)s$；若把 $T$ 也理解成某个"提升"，则等式右边将出现两个不同的对象，
%  且 $\theta(f)$ 无从出现。此处按原文排版，未改。]

因此，对应 $(\theta,s)\mapsto D(p^{-1}\theta)(s)$ 定义了 $p_*(\mathcal L)$ 在 $S$ 上的一个联络 $\nabla(D)$。

热算子 $D$ 的象征是 $\cO_{\cX}$-线性映射
\[
 \rho_D\colon p^*(T_S)\to S^2\mathcal T_{\cX/S},
\]
它由
\[
 \rho_D=\sigma^2\circ D
\]
给出，其中 $\sigma^2\colon\mathcal D_{\cX}^2(\mathcal L)\to S^2T_{\cX}$ 是象征映射。

让我们取 $S$ 的一个 $\C$-值点 $0$。令 $X=p^{-1}(0)$，$L=\mathcal L|_X$。元素 $\theta\in T_0S$ 定义了一个态射 $\theta\colon\operatorname{Spec}\C[\eps]\to S$，由基变换我们得到 $\operatorname{Spec}\C[\eps]$ 上的一对 $(X_\eps,L_\eps)$，它是 $(X,L)$ 的一阶形变。于是 $w_\theta=\rho_D(\theta)\in H^0(X,S^2T_X)$，而对任意 $s\in H^0(X,L)$，$-\beta(\delta_s(w_\theta))$ 是一阶形变 $(X_\eps,L_\eps)$ 的小平–Spencer 类。

反过来，假定给定了 $\cO_{\cX}$-线性映射 $\rho\colon p^*(T_S)\to S^2\mathcal T_{\cX/S}$，考虑能否找到 $\mathcal L$ 在 $S$ 上的热算子 $D$ 使 $\rho=\rho_D$。

\begin{lemma}[5.29]（[GJ]）\quad 若下述两个条件得到满足，则我们可以局部地找到一个热算子。
\begin{enumerate}
\item $R^1p_*\cO_{\cX}$ 与 $R^1p_*T_{\cX/S}$ 在 $S$ 上局部自由，且在 $S$ 的任一 $\C$-值点 $0$ 处，其纤维为 $H^1(X,\cO_X)$ 与 $H^1(X,T_X)$，其中 $X=p^{-1}(0)$。
\item 对 $S$ 的任一 $\C$-值点 $0$ 与任一元素 $\theta\in T_0S$，$(X_\eps,L_\eps)$ 的小平–Spencer 类由 $-\beta(\delta_s(\rho(\theta)))$ 给出。
\end{enumerate}
\end{lemma}

% [原文 p.125]

对亏格 $g$ 的光滑曲线上 $G$-丛的模空间，我们可以应用引理 5.29，从而在 $\mathcal M_g$ 中局部地总可以找到 $\mathcal L^{\otimes\ell}$ 上的热算子，使 $\rho_D=\tau$，其中 $\tau$ 由下文 (5.60) 定义。为了把 Hitchin 联络与我们在 §5.1 中引入的联络作比较，更好的办法是利用能量–动量张量（energy-momentum tensor）直接构造热算子。

我们仍令 $G$ 为如上所述的单连通单代数群，$E$ 为光滑射影曲线 $C$ 上的半稳定 $G$-丛。半稳定 $G$-丛在 $C$ 上的粗模空间 $M_G(C)$ 在对应 $E$ 的点 $[E]$ 处的切空间由 $H^1(C,\operatorname{Ad}(E))$ 给出。此外还存在自然的杯积配对
\[
 H^1(C,T_C)\otimes H^0(C,\operatorname{Ad}(E)\otimes\Omega^1_C)\to H^1(C,\operatorname{Ad}(E)).
\]
由于 Cartan–Killing 型给出 $\operatorname{Ad}(E)$ 与其对偶之间的同构，由 Serre 对偶，该配对诱导出同态
\begin{equation}
 \tau_{[E]}\colon H^1(C,T_C)\to S^2H^1(C,\operatorname{Ad}(E))=(S^2T_{M_G(C)})_{[E]}. \tag{5.56}
\end{equation}

让我们选取曲线 $C$ 上的点 $P$ 与以 $P$ 为中心在 $C$ 上的坐标 $\xi$，并考虑带坐标的带点曲线 $(C,Q,\xi)$。Cartan–Killing 型定义一个配对
\begin{equation}
 \operatorname{tr}\colon\Bigl(\mathfrak g\otimes\C((\xi))\frac{d}{d\xi}\Bigr)\otimes_{\C((\xi))}\Bigl(\mathfrak g\otimes\C((\xi))\Bigr)\to\C((\xi))\frac{d}{d\xi}. \tag{5.57}
\end{equation}

% [Tutor Note: 这里先取点 $P$、后又写 $(C,Q,\xi)$，$Q$ 应为 $P$（原文笔误）；
%  本节此后一直用 $(C,P,\xi)$。按原文排版，未改。]

对 $\ell\in\C((\xi))\frac{d}{d\xi}$，令 $\theta(\ell)\in H^1(C,T_C)$ 是由 $\ell$ 定义的元素（见 §4.2.1）。则 $\tau_{[E]}(\theta(\ell))\in S^2H^1(C,\operatorname{Ad}(E))$ 可以看作 $H^0(C,\operatorname{Ad}(E)\otimes\Omega^1_C)$ 上的一个二次型。它可由下述公式给出
\begin{equation}
 \tau_{[E]}(\theta(\ell))(\phi)=\operatorname*{Res}_{\xi=0}\operatorname{tr}\bigl(\bar\phi\otimes\langle\ell,\bar\phi\rangle\bigr) \tag{5.58}
\end{equation}
其中 $\phi$ 取遍 $H^0(C,\operatorname{Ad}(E)\otimes\Omega^1_C)$，$\bar\phi\in\mathfrak g\otimes\C((\xi))\frac{d}{d\xi}$ 由在 $P$ 处的洛朗展开得到，而 $\langle\ell,\bar\phi\rangle\in\mathfrak g\otimes\C((\xi))$ 是缩并。

现在让 $E$ 在模空间 $M_G(C)$ 上变动，由 (5.56) 我们得到同态
\begin{equation}
 \tau_C\colon H^1(C,T_C)\to H^0(M_G(C),S^2T_{M_G(C)}). \tag{5.59}
\end{equation}
此外，让曲线 $C$ 在 $\mathcal M_g$ 中变动，我们得到 $\cO_{\mathcal M_G}$-模同态
\begin{equation}
 \tau\colon p^*T_{\mathcal M_g}\to S^2T_{\mathcal M_G/\mathcal M_g}, \tag{5.60}
\end{equation}
其中 $p\colon\mathcal M_G\to\mathcal M_g$ 是自然态射。又令
\begin{equation}
 w_s=\delta_s\circ\tau_C\colon H^1(C,T_C)\to\HH^1(d_1s), \tag{5.61}
\end{equation}
其中 $\delta_s$ 由 (5.53) 定义。


注意，$\theta\in H^1(C,T_C)$ 定义了曲线 $C$ 的一个一阶形变，因此它定义了态射 $\theta_\eps\colon\operatorname{Spec}\C[\eps]\to\mathcal M_g$。令 $(\mathcal M_G(C)_\eps,C_\eps,\mathcal L^{\otimes\ell}_\eps)$ 是 $(\mathcal M_G(C),\mathcal L^{\otimes\ell})$ 沿态射 $\theta_\eps$ 的拉回，其中 $C\to\mathcal M_g$ 是亏格 $g$ 的光滑曲线的模栈上的通用曲线，而 $\mathcal M_G(C)$ 是亏格 $g$ 的曲线 $C$ 上半稳定 $G$-丛的模栈。

\begin{lemma}[5.30]（[Hi]，[L]）\quad 由 $-w_s(\theta)/(2\ell+2g^*)$ 在 $H^1(\mathcal M_G(C),\mathcal D_{\mathcal M_G(C)}^1(\mathcal L^{\otimes\ell}))$ 中的投影所定义的 $(\mathcal M_G(C),\mathcal L^{\otimes\ell})$ 的形变，同构于 $(\mathcal M_G(C)_\eps,\mathcal L^{\otimes\ell}_\eps)$，其中 $g^*$ 是与 $G$ 相伴的李代数 $\mathfrak g$ 的对偶 Coxeter 数。
\end{lemma}

% [Tutor Note: 分母 $2\ell+2g^*=2(\ell+g^*)$ 是 Hitchin 联络的"反常移位"：
%  权 $\ell$ 被替换为 $\ell+g^*$（$g^*$ 为对偶 Coxeter 数，例如 $A_n$ 时 $g^*=n+1$）。
%  这一移位是 $2\ell+2g^*$ 而非 $\ell+g^*$，因为这里的 $\ell$ 是线丛幂次（二次微分
%  的阶数），规范化后得到的因子正是 $\frac{1}{\ell+g^*}$。若读者按 $\frac{1}{2\ell}$
%  去核对，会出现因子 2 的偏差。]

% [原文 p.126]

现在让我们构造 $\mathcal L^{\otimes\ell}$ 上的射影热算子 $D$，使 $\rho_D=\tau$，其中 $\tau$ 由 (5.60) 定义。首先注意存在正合列
\[
 0\to\operatorname{Ad}(E)\to\operatorname{Ad}(E)(mQ)\to\oplus_{i=1}^m\mathfrak g\otimes\xi^{-i}\to0.
\]
由上同调正合列我们得到同态
\[
 \oplus_{i=1}^m\mathfrak g\otimes\xi^{-i}\to H^1(C,\operatorname{Ad}(E)),
\]
它诱导出同态
\begin{equation}
 \psi\colon\mathfrak g\otimes_\C\C((\xi))\to\mathfrak g\otimes_\C\C[\xi^{-1}]\xi^{-1}\to H^1(C,\operatorname{Ad}(E)). \tag{5.62}
\end{equation}

同态 $\psi$ 依赖于曲线 $C$ 的点 $P$ 的选取以及 $P$ 处的形式坐标 $\xi$。令 $\mathcal C\to\mathcal M_g$ 是亏格 $g$ 的光滑曲线的通用族。于是我们总能找到对应于曲线 $C$ 的点 $[C]\in M_g$ 的开仿射邻域 $S$ 以及通用族 $\pi\colon\mathcal C_S\to S$，使得该族有截面 $s_S$ 以及沿 $s_S$ 的坐标 $\xi$。由于 $H^1(C,\operatorname{Ad}(E))$ 是模栈 $\mathcal M_G(C)$ 在 $[E]$ 处的切空间，通过让 $C$ 在 $S$ 中变动、$E$ 变动，我们得到从 $\mathfrak g\otimes_\C\C((\xi))$ 到 $S$ 上半稳定 $G$-丛的通用模栈 $\mathcal M_{G,\pi}$ 上的全纯向量场空间的一个同态：
\[
 \mathfrak g\otimes_\C\C((\xi))\to H^0(\mathcal M_{G,\pi},T_{\mathcal M_{G,\pi}}).
\]
该同态可以延拓为来自 $\hat{\mathfrak g}$ 的同态：
\[
 \hat{\mathfrak g}\to\mathfrak g\otimes_\C\C((\xi))\to H^0(\mathcal M_{G,\pi},T_{\mathcal M_{G,\pi}}).
\]


此外，该同态还可延拓为从 $\hat{\mathfrak g}$ 的泛包络代数到 $\mathcal M_{G,\pi}$ 上微分算子环的环同态：
\begin{equation}
 U(\hat{\mathfrak g})^{\mathrm{opp}}\to H^0\bigl(\mathcal M_{G,\pi},\mathcal D_{\mathcal M_{G,\pi}}(\cO_{\mathcal M_{G,\pi}})\bigr), \tag{5.63}
\end{equation}
其中 $\mathcal D_{\mathcal M_{G,\pi}}(\cO_{\mathcal M_{G,\pi}})$ 是 $\cO_{\mathcal M_{G,\pi}}$ 的微分算子层。这里我们采用不同的约定，即向量场从右边作用，以便与我们 §3.1 中给出的共形块的构造相容。

令
\[
 T(z)=\sum_{n\in\Z}L_n t^{-n-2}
\]
是仿射李代数的能量–动量张量（见定义 2.6）。对 $n\ne0$，由 (2.12) 我们有
\[
 L_n=\frac{1}{2(g^*+\ell)}\sum_{m\in\Z}\sum_{a=1}^{\dim\mathfrak g}J^a(m)J^a(n-m),
\]
而由 (2.13) 我们有
\[
 L_0=\frac{1}{2(g^*+\ell)}\sum_{a=1}^{\dim\mathfrak g}J^aJ^a+\frac{1}{(g^*+\ell)}\sum_{m=1}^{\infty}J^a(-m)J^a(m),
\]
其中 $\{J^1,\dots,J^{\dim\mathfrak g}\}$ 是 $\mathfrak g$ 关于归一化 Cartan–Killing 型的正交基。尽管 $L_n$ 不属于 $U(\hat{\mathfrak g})$，但由 (5.62) 中 $\psi$ 的定义，$J^a(k)$（$k\ge0$）定义了零向量场，于是 $L_n$ 定义了 $\cO_{\mathcal M_{G,\pi}}$ 上的一个二阶微分算子。而且若 $n>0$，该微分算子为零。对 $\ell=\ell(\xi)\frac{d}{d\xi}\in\C((\xi))\frac{d}{d\xi}$，在表达式 $T[\ell]=\operatorname*{Res}_{z=0}(\ell(z)T(z)dz)=\sum l_n L_n$

% [Tutor Note: $T(z)=\sum L_n t^{-n-2}$ 与 $T[\ell]=\operatorname*{Res}(\ell(z)T(z)dz)$ 中的
%  变量应为同一个 $z$，原文把 $T(z)$ 的展开写成 $t^{-n-2}$（笔误）；按原文排版，未改。
%  另外末式系数原文印作 $l_n$（直立体 l），应为 $\ell_n$，即 $\ell(z)=\sum\ell_n z^{n+1}$
%  的系数。$L_0$ 的第二项按重复指标 $a$ 求和，原文未写 $\sum_a$。]

中只有有限多个 $n<0$ 的 $L_n$。因此 $T[\ell]$ 定义了二阶微分算子，于是我们有同态
\[
 \C((\xi))\frac{d}{d\xi}\to H^0\bigl(\mathcal M_{G,\pi},\mathcal D_{\mathcal M_{G,\pi}}^2(\cO_{\mathcal M_{G,\pi}})\bigr).
\]
利用族 $\mathcal C_S\to S$ 的沿 $s_S$ 带坐标 $\xi$ 的截面 $s_S$，我们可以证明下图是交换的：
\[
 \begin{array}{ccc}
 \C((\xi))\frac{d}{d\xi} & \longrightarrow & H^0\bigl(\mathcal M_{G,\pi},\mathcal D_{\mathcal M_{G,\pi}}^2(\cO_{\mathcal M_{G,\pi}})\bigr)\\[4mm]
 & \theta_S\searrow\qquad\nearrow D_S & \\[4mm]
 & H^0(S,T_S) &
 \end{array}
\]
其中同态
\[
 \theta_S\colon\C((\xi))\frac{d}{d\xi}\to H^0(S,T_S)
\]
在 §4.2 中讨论。这意味着存在 $\cO_{\mathcal M_{G,\pi}}$-线性同态
\[
 D\colon p^*T_S\to\mathcal D_{\mathcal M_{G,\pi}}^2(\cO_{\mathcal M_{G,\pi}}),
\]
其中 $p\colon\mathcal M_{G,\pi}\to S$ 是自然态射。我们可以证明它是 $\cO_{\mathcal M_{G,\pi}}$ 上的一个热算子。

% [原文 p.127]

严格说来，上述论证并不正确，因为我们需要用概形论代替栈。此外，粗模概形 $\underline M_G(C)$ 一般有奇点，因此我们需要像 [Hi] 与 [L] 中那样，选取对应稳定 $G$-丛的光滑开部分。但在下文中我们将用这些不精确的论证来说明 Hitchin 联络构造的基本想法。详细的论证可在 [Hi]、[L] 与 [GJ] 中找到。


现在让我们用与上面类似的方法说明可以在 $\mathcal M_{G,\pi}$ 上定义 $\mathcal L^{\otimes\ell}$ 上的热算子。注意我们的线丛 $\mathcal L$ 最初定义在与特征标 $\chi^{-1}\colon\widehat{L_\pi^+G}\to\mathbb G_{m,S}$ 相伴的族 $\pi\colon\mathcal C_S\to S$ 的栈 $\mathcal Q_\pi=\widehat{L_\pi G}/\widehat{L_\pi^+G}$ 上。$S$ 上的通用模栈 $p\colon\mathcal M_{G,\pi}\to S$ 作为商栈 $L_{\mathfrak F}G\backslash\mathcal Q_\pi$ 得到，其中 $\mathfrak F=(\pi\colon\mathcal C_S\to S;s_S;\xi)$，而线丛 $\mathcal L$ 下降到 $\mathcal M_{G,\pi}$ 上的一个线丛，我们仍用同一记号 $\mathcal L$ 表示它。因此 $\mathcal M_{G,\pi}$ 上 $\mathcal L^{\otimes\ell}$ 的一个局部全纯截面，对应于 $\widehat{L_\pi G}$ 上满足
\[
 f(gh)=\chi(h)^\ell f(g),\qquad \forall h\in\widehat{L_\pi^+G}
\]
的 $L_{\mathfrak F}G$-不变的局部全纯截面 $f(g)$。（严格说来，如 [L] 中那样，我们需要取 $\mathcal Q_\pi\to S$ 的拟截面 $\sigma\colon N\to\mathcal Q_\pi$。）因此 $\mathcal M_{G,\pi}$ 上向量场的右作用保持该性质，于是与上述论证类似，我们有同态
\[
 U(\mathfrak g)^{\mathrm{opp}}\to H^0\bigl((\mathcal M_{G,\pi},\mathcal D_{\mathcal M_{G,\pi}}(\mathcal L^{\otimes\ell})\bigr).
\]

% [Tutor Note: 上式两处按原文排版：(i) 左边作 $U(\mathfrak g)$ 而非 $U(\hat{\mathfrak g})$，
%  与 (5.63) 的 $U(\hat{\mathfrak g})^{\mathrm{opp}}$ 不一致；(ii) $H^0($ 之后多了一个
%  左括号 "$($"。数学上应为 $U(\hat{\mathfrak g})^{\mathrm{opp}}\to
%  H^0(\mathcal M_{G,\pi},\cD_{\mathcal M_{G,\pi}}(\mathcal L^{\otimes\ell}))$。]

此外，对任意 $\ell\in\C((\xi))\frac{d}{d\xi}$，若 $\theta_S(\ell)=\theta\in H^0(S,T_S)$，我们可以把 $T[\ell]$ 的像定义为 $\mathcal L^{\otimes\ell}$ 上的二阶微分算子。若 $\theta_S(\ell)=\theta_S(\ell')$，则 $T[\ell]$ 与 $T[\ell']$ 定义同一个二阶微分算子。因此我们可以定义一个 $\cO_{\mathcal M_{G,\pi}}$-线性同态
\[
 D\colon p^*T_S\to\mathcal D_{\mathcal M_{G,\pi}}^2(\mathcal L^{\otimes\ell})
\]

% [原文 p.128]

它就是所要的热算子。按我们的构造，容易证明下述图是交换的。
\[
 \begin{array}{c}
 p^*T_S\;\xrightarrow{\;D\;}\;\mathcal D_{\mathcal M_{G,\pi}}^2(\mathcal L^{\otimes\ell})\;\xrightarrow{\ \text{象征}\ }\;S^2T_{\mathcal M_{G,\pi}/S}\\[3mm]
 p^*T_S\;\xrightarrow{\;D\;}\;\mathcal D_{\mathcal M_{G,\pi}}^2(\cO_{\mathcal M_{G,\pi}})\;\xrightarrow{\ \text{象征}\ }\;S^2T_{\mathcal M_{G,\pi}/S}
 \end{array}
\]

% [Tutor Note: 原文的这个"交换图"其实是一个叉形图：$\mathcal D^2(\mathcal L^{\otimes\ell})$
%  与 $\mathcal D^2(\cO)$ 上方的"symbol"箭头都指向同一个 $S^2T_{\mathcal M_{G,\pi}/S}$。
%  中文版为排版方便画成两行阶梯形，两条"象征"箭头指向同一对象，交换性内容不变。]

让我们对 $\ell\in\C((\xi))\frac{d}{d\xi}$ 计算象征 $D(\theta(\ell))$。取 $\phi\in H^0(C,\operatorname{Ad}(E))\otimes\Omega^1_C$，其中 $C=\pi^{-1}(x)$，$x\in S$，$E$ 是 $C$ 上的半稳定 $G$-丛；令 $\bar\phi\in\mathfrak g\otimes\C((\xi))\frac{d}{d\xi}$ 由在 $s_S(x)$ 处关于坐标 $\xi$ 的洛朗展开得到。定义
\begin{equation}
 \langle\phi\otimes\phi,L_n\rangle=\frac{1}{2(g^*+\ell)}\sum_{m\in\Z}\sum_{a=1}^{\dim\mathfrak g}\langle\phi,\psi(J^a(-m))\rangle\cdot\langle\phi,\psi(J^a(-m+n))\rangle \tag{5.64}
\end{equation}
其中 $\psi$ 由 (5.62) 定义，而 $\langle\ ,\ \rangle\colon H^0(C,\operatorname{Ad}(E)\otimes\Omega^1_C)\times H^1(C,\operatorname{Ad}(E))\to\C$ 是对偶配对。在我们的情形它由
\[
 \langle\phi,\psi(X(m))\rangle=\operatorname*{Res}_{\xi=0}\operatorname{tr}\langle\bar\phi,\otimes X(m)\rangle
\]
给出，其中 $\operatorname{tr}$ 由 (5.57) 给出。

% [Tutor Note: 上式右侧按原文排版作 $\operatorname{tr}\langle\bar\phi,\otimes X(m)\rangle$，
%  括号形状与原意都不对：应为 $\operatorname{tr}(\bar\phi\otimes X(m))$，即先作缩并
%  再取 $\operatorname{tr}$。原文的 $\langle$ 与多出的逗号是排印之误。]

由于
\[
 \langle\phi,\psi(J^a(-m))\rangle\cdot\langle\phi,\psi(J^a(-m+n))\rangle\ne0
\]
只对有限多个 $m$ 成立，(5.64) 是良定义的。此外，对固定的 $\phi$，只有有限多个 $n$ 使 $\langle\phi\otimes\phi,L_n\rangle\ne0$。因此对
\[
 \ell=\sum_{n=-m_0}^{\infty}l_n\xi^{n+1}\frac{d}{d\xi}\in\C((\xi))\frac{d}{d\xi},
\]
配对
\[
 \langle\phi\otimes\phi,T[\ell]\rangle=\sum_{n=-m_0}^{\infty}l_n\langle\phi\otimes\phi,L_n\rangle
\]
是良定义的。现在 $\bar\phi$ 具有形状
\[
 \sum_{n\in\Z}X_n\xi^{-n-1}d\xi,\qquad X_n\in\mathfrak g,
\]
其中当 $n$ 充分大时 $X_n=0$。于是我们有
\begin{align}
 \Bigl\langle\phi\otimes\phi,T\Bigl[\xi^{n+1}\frac{d}{d\xi}\Bigr]\Bigr\rangle &=\langle\phi\otimes\phi,L_n\rangle \nonumber\\
 &=\sum_{m\in\Z}\sum_a^{\dim\mathfrak g}\langle\bar\phi,J^a(-m)\rangle\cdot\langle\bar\phi,J^a(n-m)\rangle \nonumber\\
 &=\sum_{m\in\Z}\sum_a^{\dim\mathfrak g}(X_m,J^a)\cdot(X_{n-m},J^a) \nonumber\\
 &=\sum_{m\in\Z}(X_m,X_{n-m}), \nonumber
\end{align}

% [Tutor Note: 上述四行等式原文没有编号，故一律 \nonumber。第二行原文的尖括号
%  残缺（第一个因子缺右尖括号，第二个因子上方有一道横线），按上一行的同一构造
%  补全为 $\langle\bar\phi,J^a(-m)\rangle\cdot\langle\bar\phi,J^a(n-m)\rangle$；
%  这是本页唯一一处未逐字符照排的公式，特此说明。]

% [原文 p.129]

其中 $(\ ,\ )$ 是归一化 Cartan–Killing 型。注意右端的最后一个和是有限和。这给出了象征 $D\bigl(\theta_S(\xi^{n+1}\frac{d}{d\xi})\bigr)$ 在 $\phi\in H^0(C,\operatorname{Ad}(E)\otimes\Omega^1_C)$ 处的值。另一方面，由 (5.58) 我们有等式
\[
 \begin{aligned}
 \tau_{[E]}\Bigl(\theta\Bigl(\xi^{n+1}\frac{d}{d\xi}\Bigr)\Bigr)(\phi)
 &=\operatorname*{Res}_{\xi=0}\operatorname{tr}\Bigl(\bar\phi\otimes\Bigl\langle\xi^{n+1}\frac{d}{d\xi},\bar\phi\Bigr\rangle\Bigr)\\
 &=\sum_m(X_m,X_{n-m}).
 \end{aligned}
\]

于是我们得到如下重要的事实。

\begin{theorem}[5.31]
上述 $\mathcal L^{\otimes\ell}$ 在 $S$ 上的热方程的象征映射等于 $\tau|_S$，其中 $\tau$ 由 (5.60) 给出。
\end{theorem}

% [Tutor Note: 定理 5.31 原文作"heat equation"（热方程），应为"heat operator"（热算子）；
%  同句"where $\tau$ given in (5.60)"缺系动词，照原文译作"其中 $\tau$ 由 (5.60) 给出"。
%  另外下文"constructed form the energy-momentum tensor"中的"form"应为"from"。这些
%  都是排印错误，不影响数学内容。]

热算子的上述构造在 $\mathcal M_g$ 中是局部的。但我们可以证明，这样构造出的局部热算子定义了 $\mathcal L^{\otimes\ell}$ 在 $\mathcal M_G$ 上的射影热算子。因此我们的射影丛具有射影联络，它就是 Hitchin 联络。

由于射影热算子是由能量–动量张量构造出来的，由同构 (5.49) 以及我们对共形块层上联络的定义（见 (5.5)），我们推断出下述定理。

\begin{theorem}[5.32]
在同构 (5.50) 之下，Hitchin 联络与我们的联络重合。
\end{theorem}
